\section{Derivación de la parametrización reducida de Grover}
\label{ap:grover_reducido}

Este apéndice resume la derivación de la formulación reducida usada como punto de
comparación para Grover en la evaluación experimental. El objetivo es mostrar, de
forma compacta, por qué bajo las hipótesis habituales del algoritmo la dinámica ideal
puede describirse con solo dos parámetros.

Sea $W\subseteq\{0,\dots,N-1\}$ el conjunto de estados marcados por el oráculo, con
$M=|W|$. Cuando Grover parte de la superposición uniforme y el oráculo introduce únicamente
un cambio de fase sobre los estados de $W$, el espacio relevante de evolución puede
describirse mediante los vectores normalizados
\[
\ket{w}=\frac{1}{\sqrt{M}}\sum_{x\in W}\ket{x},
\qquad
\ket{r}=\frac{1}{\sqrt{N-M}}\sum_{x\notin W}\ket{x}.
\]
El estado inicial uniforme se descompone entonces como
\[
\ket{s}
=
\frac{1}{\sqrt{N}}\sum_{x=0}^{N-1}\ket{x}
=
\sqrt{\frac{M}{N}}\ket{w}+\sqrt{\frac{N-M}{N}}\ket{r}.
\]

La iteración de Grover está dada por el producto del oráculo $O$ y el operador de difusión
$D=2\ket{s}\!\bra{s}-I$. El oráculo actúa como
\[
O\ket{w}=-\ket{w},
\qquad
O\ket{r}=\ket{r},
\]
mientras que el operador de difusión refleja el estado respecto a la dirección $\ket{s}$.
Como ambos operadores preservan el subespacio generado por $\{\ket{w},\ket{r}\}$, cualquier
estado de la evolución ideal puede escribirse como
\[
\ket{\psi_t}=\alpha_t\ket{w}+\beta_t\ket{r}.
\]

Si se desea recuperar las amplitudes a nivel de estado base, basta definir
\[
u_t=\frac{\alpha_t}{\sqrt{M}},
\qquad
v_t=\frac{\beta_t}{\sqrt{N-M}},
\]
de manera que cada estado marcado tiene amplitud $u_t$ y cada estado no marcado tiene
amplitud $v_t$. Tras aplicar el oráculo, las amplitudes pasan a ser $-u_t$ sobre $W$ y
$v_t$ fuera de $W$. El promedio de amplitud en ese punto es
\[
\mu_t=\frac{-Mu_t+(N-M)v_t}{N}.
\]

El operador de difusión actúa componente a componente mediante la transformación
$x\mapsto 2\mu_t-x$. Por tanto, para un estado marcado se obtiene
\[
u_{t+1}=2\mu_t-(-u_t)=2\mu_t+u_t
=
\frac{N-2M}{N}u_t+\frac{2(N-M)}{N}v_t,
\]
y para un estado no marcado
\[
v_{t+1}=2\mu_t-v_t
=
-\frac{2M}{N}u_t+\frac{N-2M}{N}v_t.
\]
En forma matricial,
\[
\begin{bmatrix}
u_{t+1} \\
v_{t+1}
\end{bmatrix}
=
\begin{bmatrix}
\frac{N-2M}{N} & \frac{2(N-M)}{N} \\
-\frac{2M}{N} & \frac{N-2M}{N}
\end{bmatrix}
\begin{bmatrix}
u_t \\
v_t
\end{bmatrix}.
\]

La evolución completa queda así reducida a una recurrencia lineal de dimensión dos. En
consecuencia, si el objetivo es seguir la dinámica ideal de Grover bajo estas hipótesis, cada
iteración puede resolverse actualizando solo los parámetros $u_t$ y $v_t$, en lugar de
aplicar el oráculo y la difusión sobre las $N$ amplitudes del vector de estado. La
reconstrucción explícita del vector completo solo es necesaria cuando se requiere materializar
las amplitudes individuales:
\[
a_x^{(t)}=
\begin{cases}
u_t, & x\in W,\\
v_t, & x\notin W.
\end{cases}
\]

Esta reducción depende de que la preparación inicial, el oráculo y la dinámica posterior
preserven la igualdad de amplitud dentro de las clases marcada y no marcada. Por ello, se
trata de una optimización específica para Grover y no de un sustituto de un esquema general de
simulación o almacenamiento.
